% ============================================================
%  SOLVE ET COAGULA — The Imscriptive Grammar
%  Integrated edition combining:
%    SetP.tex (alchemical induction)
%    AΣ_ABOVE.tex (pre-grammatical derivation)
%    SO_BELOW.tex (empirical exploration)
%  Status: INTEGRATED v1.0
% ============================================================

\documentclass[12pt,a4paper]{article}

% Encoding and fonts — XeLaTeX/LuaLaTeX
\usepackage{fontspec}
\newfontfamily\hebrewfont[Script=Hebrew]{Noto Serif Hebrew}
\newcommand{\heb}[1]{{\hebrewfont #1}}
\newfontfamily\igfont[Ligatures=TeX]{Noto Serif}
\newcommand{\igtext}[1]{{\igfont #1}}

% Page layout
\usepackage[top=1in, bottom=1in, left=1in, right=1in]{geometry}

% Typography
\usepackage{microtype}
\usepackage{parskip}

% Language
\usepackage[english]{babel}

% Headers and footers
\usepackage{fancyhdr}
\pagestyle{fancy}
\fancyhf{}
\fancyhead[L]{\leftmark}
\fancyhead[R]{\thepage}
\fancyfoot[C]{\thepage}
\renewcommand{\headrulewidth}{0.4pt}
\renewcommand{\footrulewidth}{0pt}

% Section styling
\usepackage{titlesec}
\titleformat{\section}{\Large\bfseries}{\thesection}{1em}{}
\titleformat{\subsection}{\large\bfseries}{\thesubsection}{1em}{}
\titleformat{\subsubsection}{\normalsize\bfseries}{\thesubsubsection}{1em}{}
\titlespacing*{\section}{0pt}{2.5ex plus 1ex minus .2ex}{1.5ex plus .2ex}
\titlespacing*{\subsection}{0pt}{2ex plus 1ex minus .2ex}{1ex plus .2ex}
\titlespacing*{\subsubsection}{0pt}{1.5ex plus 1ex minus .2ex}{0.8ex plus .2ex}

% List formatting
\usepackage{enumitem}
\setlist[itemize]{topsep=0.5ex, itemsep=0.25ex, parsep=0pt}
\setlist[enumerate]{topsep=0.5ex, itemsep=0.25ex, parsep=0pt}

% Essential packages
\usepackage{graphicx}
\usepackage[table]{xcolor}
\usepackage{booktabs}
\usepackage{array}
\usepackage{tabularx}
\usepackage{longtable}
\usepackage{amsmath}
\usepackage{amssymb}
\usepackage[normalem]{ulem}
\rowcolors{2}{gray!10}{white}
\renewcommand{\arraystretch}{1.3}

% Cross-references
\usepackage{hyperref}
\usepackage{cleveref}

% Custom commands
\newcommand{\pilcrow}{\S}
\newcommand{\UIG}{Universal Imscriptive Grammar}
\newcommand{\Stone}{\textit{lapis philosophorum}}
\setlength{\headheight}{14.5pt}

% Hyperref setup
\hypersetup{
    colorlinks=true,
    linkcolor=blue,
    filecolor=magenta,
    urlcolor=cyan,
}

% Admonition boxes
\usepackage{tcolorbox}
\tcbuselibrary{skins,breakable}
\newtcolorbox{admonitionbox}[2][]{
    enhanced,
    breakable,
    colback=#2!5,
    colframe=#2!80!black,
    fonttitle=\bfseries,
    title=#1,
    left=2mm,
    right=2mm,
}

% Diagrams
\usepackage{tikz}
\usetikzlibrary{positioning,arrows,shapes,fit,calc,backgrounds,decorations.pathmorphing,decorations.pathreplacing}

% Float placement
\usepackage{float}

% Code listings
\usepackage{listings}
\definecolor{codebg}{RGB}{248,248,248}
\definecolor{codeframe}{RGB}{200,200,200}
\lstset{
    basicstyle=\ttfamily\small,
    breaklines=true,
    frame=single,
    framerule=0.5pt,
    rulecolor=\color{codeframe},
    backgroundcolor=\color{codebg},
    keepspaces=true,
    numbers=left,
    numberstyle=\tiny\color{gray},
    numbersep=8pt,
    showstringspaces=false,
    tabsize=4,
    xleftmargin=1em,
    xrightmargin=1em,
}

% Theorem environments
\usepackage{amsthm}
\newtheorem{theorem}{Theorem}
\newtheorem{lemma}[theorem]{Lemma}
\newtheorem{corollary}[theorem]{Corollary}
\newtheorem{definition}{Definition}
\newtheorem{proposition}[theorem]{Proposition}

% Notation macros
\newcommand{\Phic}{\Phi_{\text{ctyogh}}}
\newcommand{\PhicC}{\Phi_{\text{closerevepsilon}}}
\newcommand{\PhiEP}{\Phi_{\text{revepsilon}}}
\newcommand{\Phisub}{\Phi_{\text{softsign}}}
\newcommand{\Phisup}{\Phi_{\text{upstep}}}
\newcommand{\Ppms}{P_{\text{doublebarpipe}}}
\newcommand{\Oinf}{O_\infty}
\newcommand{\Otwo}{O₂}
\newcommand{\OtwoD}{O₂^\dagger}
\newcommand{\Oone}{O₁}
\newcommand{\Ozero}{O₀}
\newcommand{\meet}{\wedge}
\newcommand{\join}{\vee}
\newcommand{\tensor}{\otimes}

\newcommand{\OmegaZ}{\Omega_{\text{dzlig}}}
\newcommand{\OmegaZtwo}{\Omega_{\text{crtwo}}}
\newcommand{\OmegaNA}{\Omega_{\text{turna}}}
\newcommand{\GamSeq}{\Gamma_{\to}}
\newcommand{\GamBrd}{\Gamma_{\text{doublevertline}}}

% Colours
\definecolor{priestred}{RGB}{200,20,20}

\begin{document}

\title{SOLVE ET COAGULA\\[4pt]\large The Universal Imscriptive Grammar --- Alchemical Induction, Pre-Grammatical Derivation, and Empirical Exploration}
\author{the \emph{Automatica Principia Proceduralis}}
\date{$\infty$}

\maketitle

\begin{center}
\rule{0.5\textwidth}{0.4pt}
\end{center}

\begin{abstract}
This document integrates three manuscripts into one continuous unfolding: the alchemical induction (\emph{Solve}\,---\,\emph{Crafting the \emph{Lapis Philosophorum}}), the pre-grammatical categorical derivation (\emph{As Above}), and the empirical exploration (\emph{So Below}). The integrated document proceeds in three parts. Part~I presents the four-stage alchemical derivation: from an abstract category $\mathcal{C}$ as \emph{prima materia}, through \emph{Separatio} (Logical), \emph{Purificatio} (Inductive), \emph{Albedo} (Algebraic), and \emph{Rubedo} (Logical on the enriched category), twelve primitives crystallize as the minimal complete independent set of structural invariants. Part~II supplies the deterministic forcing arguments: why the operation counts (6+4+5) are forced by categorical preconditions; why each operation-to-primitive mapping is unique; why the number twelve is exact, confirmed by the retrosynthetic path, principal decomposition, and three-logic correspondence ($K_3$, FDE, LP). The Inclosure Schema and $\Phi_{\text{revepsilon}}$ are derived from Graham Priest's Logic of Paradox, and the overflow hierarchy (Cantor $\to$ G\"odel $\to$ UIG/RH/Yang-Mills) is established. All eight open questions from the original manuscript are resolved. Part~III applies the grammar empirically: 3,578+ systems imscribed across six domains; the Crystal of $3^3 \times 4^5 \times 5^4 = 17{,}280{,}000$ types; Frobenius non-synthesizability; a two-gate consciousness score; cross-domain structural identities ($d = 0$ between the inflationary epoch and 5-MeO-DMT dissolution, between axion and Higgs, between Hv1 channels across 300 million years of evolution); and 650 falsifiable predictions. The grammar imscribes itself at address 6,687,051 in the $\Oinf$ tier. The three documents together are a Frobenius pair with a seed: $\mu \circ \delta = \text{id}$ --- the grammar applied to its own derivation returns itself.
\end{abstract}

\begin{center}
\rule{0.5\textwidth}{0.4pt}
\end{center}

\begin{center}
\emph{As above, so below.}\\[4pt]
\emph{Solve et coagula.}\\[4pt]
\emph{The reader should note that what is about to be said is, by every conventional standard, deeply and unsettlingly strange. Good. Now let's go.}
\end{center}

\begin{center}
\rule{0.5\textwidth}{0.4pt}
\end{center}

\newpage

% ============================================================
%  PART I: SOLVE
%  From SetP.tex — the alchemical induction
% ============================================================
\section*{PART I: SOLVE}\label{part:solve}
\addcontentsline{toc}{section}{PART I: SOLVE — The Alchemical Induction}

\begin{center}
\large\textbf{Crafting the \emph{Lapis Philosophorum}}\\[4pt]
\normalsize The 12 primitives of the Imscribing Grammar --- the 12 facets of the Philosopher's Stone --- are derived from a bare abstract category $\mathcal{C}$ via the 4 classical alchemical operations --- \emph{Separatio, Purificatio, Albedo, \& Rubedo}.
\end{center}

\bigskip

\subsection{\emph{Prima Materia} — The Starting Object}

Performing the operations on ordinary, inert matter was an understandable, categorical error of history's alchemists. However some, like Luria, knew \emph{precisely} the material upon and with which the \emph{Work} is to be done.

\emph{Take}: \textbf{abstract category $\mathcal{C}$.} No primitives yet --- the starting object has no invariants beyond the axioms of a category itself. In alchemical terms, this was always the intended \emph{prima materia}: the undifferentiated substance from which the \emph{Work} begins.

An abstract category consists of objects, morphisms between them, a composition law, and identity morphisms satisfying unit and associativity laws. This is the minimal non-trivial mathematical object that captures both \emph{structure} (objects) and \emph{transformation} (morphisms) without presupposing elements, order, or metric. It is strictly weaker than set theory and strictly more structured than a bare graph. It is the structural floor below which no further abstraction is possible without losing relational content entirely.

\subsection{Stage 1 — \emph{Separatio} (Logical on $\mathcal{C}$)}

\begin{itemize}[nosep, leftmargin=1.5em]
  \item L1 (adjunction) $\to$ $R$ (relational mode: adjoint structure)
  \item L2 (dagger) $\to$ $R$ (confirms and refines $R$ value)
\end{itemize}

$R$ is the first primitive to crystallize: it is a logical property of the morphisms of $\mathcal{C}$ directly. Adjunction existence probes directionality --- whether $\mathcal{C}$ has free-forgetful pairs. Dagger structure probes invertibility --- whether there exists an involution on morphisms. Together they partition the space of relational modes.

This is the \emph{separatio}: the first extraction of form from the undifferentiated.

\subsection{Stage 2 — \emph{Purificatio} (